% !TEX program = xelatex
% International industry résumé; translated and refined from texs/sections.tex.
\documentclass{resume}
\geometry{left=0.6in,right=0.6in,top=0.5in,bottom=0.5in}
\usepackage{needspace}
\titlespacing*{\section}{0pt}{1.1ex}{0.8ex}
\titlespacing*{\subsection}{0pt}{1.0ex}{0.4ex}
\setlist[itemize]{noitemsep,topsep=0.1em,leftmargin=1.25pc}
\titleformat{\subsection}{\normalsize\raggedright}{}{0em}{}
\newcommand{\project}[1]{\par\addvspace{0.7ex}\Needspace{4\baselineskip}\textbf{#1}\par}
\newcommand{\publication}[3]{{\small\Needspace{\baselineskip}\datedline{\textbf{#1} / \textit{#2}}{\mbox{#3}}}}

\begin{document}
\pagenumbering{gobble}
\name{Hulin Wang}
\basicInfo{\email{hwang551@asu.edu}\dotSep
  \homepage[hulinwang.me]{https://www.hulinwang.me/}\dotSep
  {\large Tempe, AZ, USA}}
\basicInfo{\linkedin[linkedin.com/in/hulin-wang-4ab350164]{https://www.linkedin.com/in/hulin-wang-4ab350164/}\dotSep
  \github[github.com/DeviRule]{https://github.com/DeviRule}}

\sectionTitle{Education}{\faiconsixbf{graduation-cap}}
\datedsubsection{\textbf{Arizona State University}, Tempe, AZ}{Aug 2022 -- Expected Dec 2027}
\normalline{PhD Candidate in Computer Science \hfill GPA: 4.0/4.0}
\normalline{SEFCOM Security Lab; Advisors: Tiffany Bao and Fish Wang}
\datedsubsection{\textbf{University of Southern California}, Los Angeles, CA}{Aug 2020 -- May 2022}
\normalline{M.S. in Computer Science \hfill GPA: 3.6/4.0}
\datedsubsection{\textbf{University of California, Davis}, Davis, CA}{Sep 2016 -- Jun 2020}
\normalline{B.S. in Computer Science \hfill GPA: 3.7/4.0}

\sectionTitle{Work \& Research Experience}{\faiconsixbf{building-user}}
\datedsubsection{\textbf{Amazon}, New York, NY}{May 2026 -- Aug 2026}
\role{Applied Scientist Intern}{}
\project{Multi-Agent Triage for Static Analysis Rule Development}
\begin{itemize}
\item Reduced validation time for new static analysis rules from 10--15 days to 2--3 days by automating triage of production findings at 90\% accuracy, averaging 4 minutes and \$0.50 per finding.
\item Built a three-stage workflow with the Strands Agent SDK to summarize findings, gather code context, and verify results; scored 14 points higher than Claude Code using the same model.
\item Supported all 47 rules from two internal analyzers across Python, Java, and JavaScript without tuning for individual rules or languages.
\item Curated a benchmark of 600+ findings from one year of production triage records, covering every rule; achieved 95\% verdict consistency across five runs.
\end{itemize}

\datedsubsection{\textbf{Arizona State University, SEFCOM Security Lab}, Tempe, AZ}{Aug 2022 -- Present}
\role{Security Researcher \& PhD Candidate}{}
\project{Root-Cause-Driven Vulnerability Repair}
\begin{itemize}
\item Designed a two-stage repair pipeline combining static and dynamic analysis to locate potential root causes, capturing defects spanning multiple code locations and indirect calls missed by static-only tools such as CodeQL.
\item Supplied root-cause candidates to a lightweight patching agent with repository search and file-editing tools, enabling repairs informed by the relevant code context.
\item Built a manually annotated patch-quality dataset covering 178 C/C++ vulnerabilities to compare repairs generated by multiple agents.
\item Improved patch quality by \textbf{20\%} over baseline methods, as validated by majority vote from three human experts.
\end{itemize}

\project{UAV Firmware Security \& Cyber-Physical Vulnerability Discovery}
\normalline{DARPA FIRE, an \$8M-funded research project}
\begin{itemize}
\item Led a sim-to-real testing pipeline for ArduPilot and a commercial derivative: fingerprinted firmware, recreated targets in PX4/Gazebo software-in-the-loop (SITL) simulation, and applied policy-guided fuzzing (PGFuzz) to parameters, commands, and physical states.
\item Reproduced two known policy-violating input sets on physical airframes, demonstrating that reduced-fidelity simulation can expose real hardware behavior without a digital twin.
\item Established ARM Cortex-M bare-metal firmware rehosting: reconstructed ELF images from flash dumps with custom linker scripts, patched architecture checks and unmodeled MMIO reads for QEMU execution, and added GDB instrumentation for live memory inspection.
\item Recovered symbols from stripped firmware by inferring the original toolchain, cross-compiling matching libraries, and applying FLIRT signature matching in IDA Pro/Ghidra.
\end{itemize}

\newpage
\normalline{\textbf{Arizona State University, SEFCOM Security Lab} \textit{(continued)}}
\project{LLM-Assisted Fuzzing}
\begin{itemize}
\item Built an LLM-assisted source patching and AFL++ fuzzing pipeline to uncover \emph{occluded bugs}: vulnerabilities masked by more readily triggered crashes in mature codebases.
\item Automated crash clustering, root-cause hints, LLM-guided edits, rebuilding, and repeated fuzzing; used targeted regression tests to check patches, reducing false positives from 90\% to 30\% and producing vulnerability reports and proof-of-concept (PoC) test cases.
\item Found 20+ distinct occluded vulnerabilities per target in projects including Objdump, FFmpeg, and GPAC, covering use-after-free and heap/stack overflows; validated findings with ASan/MSan and targeted regression tests.
\end{itemize}

\datedsubsection{\textbf{Team Shellphish / DARPA AIxCC}, Tempe, AZ}{Aug 2023 -- Aug 2025}
\role{Patching Component Lead}{}
\normalline{\textbf{5th place; \$2M team prize} in the DARPA Artificial Intelligence Cyber Challenge (AIxCC), a global competition for automated vulnerability discovery and repair in open-source and critical infrastructure software.}
\project{Vulnerability Patching Agent}
\begin{itemize}
\item Led development of a multi-agent vulnerability repair system integrating autonomous triage, patch generation, and validation across distributed services.
\item Built a feedback loop combining fuzzing, proof-of-vulnerability re-tests, and deterministic regression tests to accept patches that eliminate vulnerabilities and preserve functionality; achieved \textbf{$>$95\% patch correctness} in the final scored rounds.
\item Triaged and patched 30+ C/Java vulnerabilities in projects including curl, NGINX, and Apache ZooKeeper, covering buffer overflows, use-after-free, command injection, and path traversal.
\end{itemize}
\project{Proof-of-Concept Generation Agent}
\begin{itemize}
\item Developed a multi-agent system combining static analysis results and code context to generate test cases targeting specific vulnerabilities.
\item Built a reachability-based program slicing pipeline to focus agents on relevant code paths, using coverage-guided search to bypass blocking logic.
\item Generated PoCs in an average of 5 minutes each, triggering 10 vulnerabilities in ZooKeeper, Tika, and jsoup (10K--100K lines of code); five were missed by Jazzer after 8 hours of fuzzing.
\end{itemize}

\datedsubsection{\textbf{USC ITS Governance, Risk \& Compliance}, Los Angeles, CA}{Oct 2021 -- May 2022}
\role{Student Researcher}{}
\project{Campus IT Risk Assessment \& Vulnerability Management}
\begin{itemize}
\item Designed and deployed four risk metrics and monitoring systems for university-wide infrastructure, building scalable Python analytics to quantify security risk in real time.
\item Developed an analytical framework to evaluate and improve BigFix patch distribution efficiency; presented findings and actionable recommendations to university leadership.
\item Integrated data from SentinelOne and other sources to calculate risk indicators and automatically alert security operations teams.
\end{itemize}

\datedsubsection{\textbf{LitePoint}, San Jose, CA}{Jun 2019 -- Sep 2019}
\role{Software Engineer Intern}{}
\project{Manufacturing Data ETL \& Query Optimization}
\begin{itemize}
\item Built a Python extract, transform, and load (ETL) system to ingest raw manufacturing data into Microsoft SQL Server.
\item Automated retrieval and cleaning of 10GB of new CSV data from designated servers, loading it into the database within 2 minutes using parallel streams.
\item Improved SQL query performance by 500\% through indexing, optimizer hints, and refactoring legacy queries.
\end{itemize}

\sectionTitle{Publications}{\faiconsixbf{book}}
\publication{Root-Cause-Driven Automated Vulnerability Repair}{First author; under submission}{Apr 2026}
\publication{SeedSmith: LLM-Driven Seed Synthesis for Directed Fuzzing}{Under submission}{Apr 2026}
\publication{BuildBench: Benchmarking LLM Agents on Compiling Real-World Open-Source Software}{TMLR 2026}{Apr 2026}
\publication{AIJon: Automated Generation of Annotations for Fuzzing}{arXiv preprint}{Feb 2026}
\publication{FedNLP: A Research Platform for Federated Learning in Natural Language Processing}{NAACL Findings}{Apr 2022}

\sectionTitle{Selected Honors}{\faiconsixbf{award}}
\datedline{DARPA AIxCC, 5th place worldwide (Team Shellphish)}{Aug 2025}
\datedline{DEF CON CTF Finals, 11th place (Team Shellphish)}{2023, 2025}
\datedline{ACM CCS Distinguished Artifact Reviewer Award}{Oct 2024}
% Fellowships and scholarships retained for academic or expanded versions.
% \datedline{Arizona State University SCAI Doctoral Fellowship}{Aug 2022}
% \datedline{Arizona State University Fulton Fellowship}{Aug 2022}
% \datedline{PLDI Programming Languages Mentoring Workshop (PLMW) Scholarship}{Jun 2022}

% Keep for security research / applied scientist roles; omit for general SWE roles.
\sectionTitle{Academic Service}{\faiconsixbf{clipboard-check}}
\normalline{\textbf{Journal Reviewer:} IEEE TDSC; ACM TIST (external reviewer)}
\normalline{\textbf{Artifact Evaluation Committee:} IEEE S\&P 2026; ACM CCS 2024}

\sectionTitle{Technical Skills}{\faiconsixbf{gears}}
\normalline{\textbf{Languages:} C/C++, Python, Java, OCaml, TypeScript, SQL}
\normalline{\textbf{Security \& Reverse Engineering:} AFL++, Jazzer, PGFuzz, CodeQL, IDA Pro, Ghidra, angr, QEMU, GDB}
\normalline{\textbf{LLMs \& Agents:} OpenAI Agents SDK, LangChain, PyTorch, Hugging Face Transformers}
\normalline{\textbf{Systems \& Databases:} Docker, Kubernetes, Microsoft SQL Server, SQLite, Neo4j}
\normalline{\textbf{Embedded Simulation:} PX4/Gazebo SITL}
\normalline{\textbf{Cloud Platforms:} AWS, Azure}
\end{document}
